
\section{The case of the Lie algebra of a smooth algebraic group:
the density theorem}
\SGARefTarget{sga3:xiii:section:5}

Let \(G\) be a smooth algebraic group over the field \(k\), and let
\(\mathfrak g\) be its Lie algebra. Let \(\mathfrak h\) be a Lie
subalgebra of \(\mathfrak g\). Let \(G\) act on \(W(\mathfrak g)\)
through the adjoint representation, and consider the subscheme
\(W(\mathfrak h)\). The construction of \SGARefLink{sga3:xiii:section:1}{Section~1} leads us to introduce
\[
  N=\operatorname{Norm}_G(\mathfrak h)
   =\operatorname{Norm}_G\bigl(W(\mathfrak h)\bigr),
\]
which is an algebraic subgroup of \(G\) (not necessarily smooth), and
\[
  \mathfrak n=\operatorname{Lie}(N),
\]
together with the scheme
\[
  X=G\times^N W(\mathfrak h),
\]
the quotient of \(\widetilde X=G\times W(\mathfrak h)\) by the right
action of \(N\) given by
\[
  (g,x)\cdot n=\bigl(gn,\operatorname{Ad}(n^{-1})x\bigr).
\]
We consider the canonical morphisms
\SGARefTarget{sga3:xiii:section:5:diagram:canonical-morphisms}
\[
\begin{tikzcd}[column sep=large,row sep=large]
G\times W(\mathfrak h)
  \arrow[d,"q"']
  \arrow[dr,"\varphi"] & \\
X \arrow[r,"\psi"] & W(\mathfrak g).
\end{tikzcd}
\]

\begin{theorem}
\SGARefTarget{sga3:xiii:theorem:5:1}
With the preceding notation, suppose that \(k\) is infinite. Consider
the following conditions:
\begin{enumerate}
  \item[(i)]\SGARefTarget{sga3:xiii:theorem:5:1:condition:i}
  \(\mathfrak h\) contains a Cartan subalgebra \(\mathfrak d\) of
  \(\mathfrak g\).

  \item[(ii)]\SGARefTarget{sga3:xiii:theorem:5:1:condition:ii}
  There exists \(a\in\mathfrak h\) such that
  \(\operatorname{ad}(a)_{\mathfrak g/\mathfrak h}\) is injective.

  \item[(iii)]\SGARefTarget{sga3:xiii:theorem:5:1:condition:iii}
  The morphism
  \(\varphi:G\times W(\mathfrak h)\to W(\mathfrak g)\) is generically
  smooth.

  \item[(iv)]\SGARefTarget{sga3:xiii:theorem:5:1:condition:iv}
  \SGARefLink{sga3:xiii:theorem:5:1:condition:iii}{The preceding morphism} \(\varphi\) (or, equivalently,
  \(\psi:X\to W(\mathfrak g)\)) is dominant, and \(\mathfrak h\) has
  the same nilpotent rank as \(\mathfrak g\).

  \item[(v)]\SGARefTarget{sga3:xiii:theorem:5:1:condition:v}
  The morphism \(\psi:X\to W(\mathfrak g)\) is generically smooth and
  \(\mathfrak h=\mathfrak n\).

  \item[(vi)]\SGARefTarget{sga3:xiii:theorem:5:1:condition:vi}
  The morphism \(\psi:X\to W(\mathfrak g)\) is dominant, and
  \(\mathfrak h=\mathfrak n\).

  \item[(vii)]\SGARefTarget{sga3:xiii:theorem:5:1:condition:vii}
  The morphism \(\psi:X\to W(\mathfrak g)\) is dominant.

  \item[(viii)]\SGARefTarget{sga3:xiii:theorem:5:1:condition:viii}
  The morphism \(\psi:X\to W(\mathfrak g)\) is dominant, and \(N\) is
  smooth.

  \item[(ix)]\SGARefTarget{sga3:xiii:theorem:5:1:condition:ix}
  The morphism \(\psi:X\to W(\mathfrak g)\) is dominant,
  \(\mathfrak h=\mathfrak n\), and \(\mathfrak h\) is the Lie algebra
  of a smooth algebraic subgroup \(H\) of \(G\).

  \item[(x)]\SGARefTarget{sga3:xiii:theorem:5:1:condition:x}
  The morphism \(\psi:X\to W(\mathfrak g)\) is generically
  quasi-finite, and \(\mathfrak n=\mathfrak h\).

  \item[(xi)]\SGARefTarget{sga3:xiii:theorem:5:1:condition:xi}
  The morphism \(\psi:X\to W(\mathfrak g)\) is generically étale, and
  \(\mathfrak n=\mathfrak h\).

  \item[(xii)]\SGARefTarget{sga3:xiii:theorem:5:1:condition:xii}
  There exists a smooth algebraic subgroup \(H\) of \(G\), with Lie
  algebra \(\mathfrak h\), such that \(\mathfrak h\) contains a Cartan
  subalgebra \(\mathfrak d\) of \(\mathfrak g\).

  \item[(xiii)]\SGARefTarget{sga3:xiii:theorem:5:1:condition:xiii}
  There exists \(a\in\mathfrak h\) such that \(N\) and the transporter
  \(M_a\) of \(a\) into \(\mathfrak h\) (cf.\ \SGARefLink{sga3:xiii:section:1}{Section~1}) coincide in a
  neighborhood of \(e\), and \(\mathfrak n=\mathfrak h\).
\end{enumerate}
One then has \SGARefLink{sga3:xiii:theorem:5:1:diagram:implications}{the following diagram} of implications:
\SGARefTarget{sga3:xiii:theorem:5:1:diagram:implications}
\[
\begin{tikzcd}[column sep=1.45em,row sep=3.2em]
\SGARefLink{sga3:xiii:theorem:5:1:condition:xi}{(\mathrm{xi})}
  \arrow[rr,Leftrightarrow] &&
\SGARefLink{sga3:xiii:theorem:5:1:condition:xii}{(\mathrm{xii})}
  \arrow[r,Leftrightarrow] &
\SGARefLink{sga3:xiii:theorem:5:1:condition:xiii}{(\mathrm{xiii})}
  \arrow[r,Rightarrow] &
\SGARefLink{sga3:xiii:theorem:5:1:condition:viii}{(\mathrm{viii})}
  \arrow[r,Leftrightarrow] &
\SGARefLink{sga3:xiii:theorem:5:1:condition:ix}{(\mathrm{ix})}
  \arrow[r,Leftrightarrow] &
\SGARefLink{sga3:xiii:theorem:5:1:condition:x}{(\mathrm{x})} \\
\SGARefLink{sga3:xiii:theorem:5:1:condition:i}{(\mathrm{i})}
  \arrow[r,Leftrightarrow] &
\SGARefLink{sga3:xiii:theorem:5:1:condition:ii}{(\mathrm{ii})}
  \arrow[r,Leftrightarrow] &
\SGARefLink{sga3:xiii:theorem:5:1:condition:iii}{(\mathrm{iii})}
  \arrow[u,Leftarrow]
  \arrow[r,Leftrightarrow] &
\SGARefLink{sga3:xiii:theorem:5:1:condition:iv}{(\mathrm{iv})}
  \arrow[r,Rightarrow] &
\SGARefLink{sga3:xiii:theorem:5:1:condition:v}{(\mathrm{v})}
  \arrow[r,Rightarrow] &
\SGARefLink{sga3:xiii:theorem:5:1:condition:vi}{(\mathrm{vi})}
  \arrow[u,Leftarrow]
  \arrow[r,Rightarrow] &
\SGARefLink{sga3:xiii:theorem:5:1:condition:vii}{(\mathrm{vii})}.
\end{tikzcd}
\]
When \(k\) has characteristic zero, all the conditions under
consideration are equivalent. Finally,
\[
  \SGARefLink{sga3:xiii:theorem:5:1:condition:xi}{(\mathrm{xi})}
  \Longleftrightarrow
  [\SGARefLink{sga3:xiii:theorem:5:1:condition:i}{(\mathrm{i})}\ \text{and}\ \SGARefLink{sga3:xiii:theorem:5:1:condition:viii}{(\mathrm{viii})}]
  \Longleftrightarrow
  [\SGARefLink{sga3:xiii:theorem:5:1:condition:v}{(\mathrm{v})}\ \text{and}\ \SGARefLink{sga3:xiii:theorem:5:1:condition:viii}{(\mathrm{viii})}].
\]
\end{theorem}

Let us first record the trivial implications
\[
  \SGARefLink{sga3:xiii:theorem:5:1:condition:v}{(\mathrm{v})}\Rightarrow\SGARefLink{sga3:xiii:theorem:5:1:condition:vi}{(\mathrm{vi})}\Rightarrow\SGARefLink{sga3:xiii:theorem:5:1:condition:vii}{(\mathrm{vii})},
  \qquad
  \SGARefLink{sga3:xiii:theorem:5:1:condition:ix}{(\mathrm{ix})}\Rightarrow\SGARefLink{sga3:xiii:theorem:5:1:condition:vi}{(\mathrm{vi})},
  \qquad
  \SGARefLink{sga3:xiii:theorem:5:1:condition:xi}{(\mathrm{xi})}
  \Longleftrightarrow
  [\SGARefLink{sga3:xiii:theorem:5:1:condition:x}{(\mathrm{x})}\ \text{and}\ \SGARefLink{sga3:xiii:theorem:5:1:condition:v}{(\mathrm{v})}].
\]

We prove the equivalence of conditions \SGARefLink{sga3:xiii:theorem:5:1:condition:i}{(i)} through \SGARefLink{sga3:xiii:theorem:5:1:condition:iv}{(iv)}, and that they
imply \SGARefLink{sga3:xiii:theorem:5:1:condition:v}{(v)}. The implication \SGARefLink{sga3:xiii:theorem:5:1:condition:i}{(i)}\(\Rightarrow\)\SGARefLink{sga3:xiii:theorem:5:1:condition:ii}{(ii)} is immediate. On the
other hand, when \(k\) is algebraically closed, condition \SGARefLink{sga3:xiii:theorem:5:1:condition:iii}{(iii)} says
that there is a \(k\)-rational point of \(G\times W(\mathfrak h)\) at
which the tangent map of \(\varphi\) is surjective. It is clear that,
after translating by a suitable element of \(G(k)\), this point may be
taken to have the form \((e,a)\), with \(a\in\mathfrak h\). It follows
that when \(k\) is infinite, though not necessarily algebraically
closed, \SGARefLink{sga3:xiii:theorem:5:1:condition:iii}{this condition}---which is plainly sufficient for generic
smoothness---is also necessary.

The tangent map is readily computed. Upon identifying the tangent
space of \(W(\mathfrak h)\) at \(a\) with \(\mathfrak h\), it is the
map
\[
  (\xi,x)\longmapsto[\xi,a]+x
\]
from \(\mathfrak g\times\mathfrak h\) to \(\mathfrak g\). Its
surjectivity is equivalent to the surjectivity of
\(\operatorname{ad}(a)_{\mathfrak g/\mathfrak h}\), or, what amounts
to the same thing, to its injectivity. This proves the equivalence of
\SGARefLink{sga3:xiii:theorem:5:1:condition:ii}{(ii)} and \SGARefLink{sga3:xiii:theorem:5:1:condition:iii}{(iii)}. Moreover, \SGARefLink{sga3:xiii:theorem:5:1:condition:ii}{(ii)} plainly implies
\(\mathfrak h=\mathfrak n\), and \SGARefLink{sga3:xiii:theorem:5:1:condition:iii}{(iii)} implies that \(\psi\) is
generically smooth: if \(\varphi\) is smooth at a point \(u\), then,
since \(q\) is flat, \(\psi\) is smooth at \(q(u)\). Thus \SGARefLink{sga3:xiii:theorem:5:1:condition:ii}{(ii)} and
\SGARefLink{sga3:xiii:theorem:5:1:condition:iii}{(iii)} imply \SGARefLink{sga3:xiii:theorem:5:1:condition:v}{(v)}.

We now show that they imply \SGARefLink{sga3:xiii:theorem:5:1:condition:i}{(i)}. Since \(\psi\) is dominant and the set
of regular points of \(\mathfrak g\) is dense and open, \(\mathfrak h\)
contains regular elements of \(\mathfrak g\). Hence a sufficiently
general element \(b\) of \(\mathfrak h\) is regular in \(\mathfrak g\)
and makes \(\operatorname{ad}(b)_{\mathfrak g/\mathfrak h}\)
injective. Consequently
\(\operatorname{Nil}(b,\mathfrak g)\subset\mathfrak h\), so
\(\mathfrak h\) contains the Cartan subalgebra
\(\mathfrak d=\operatorname{Nil}(b,\mathfrak g)\). Conditions \SGARefLink{sga3:xiii:theorem:5:1:condition:i}{(i)},
\SGARefLink{sga3:xiii:theorem:5:1:condition:ii}{(ii)}, and \SGARefLink{sga3:xiii:theorem:5:1:condition:iii}{(iii)} are therefore equivalent. Finally,
\SGARefLink{sga3:xiii:theorem:5:1:condition:i}{(i)}\(\Longleftrightarrow\)\SGARefLink{sga3:xiii:theorem:5:1:condition:iv}{(iv)}: we have already observed that if
\(\mathfrak h\) contains a Cartan subalgebra, then it has the same
nilpotent rank as \(\mathfrak g\) (\SGARefLink{sga3:xiii:proposition:4:6}{4.6}), so \SGARefLink{sga3:xiii:theorem:5:1:condition:i}{(i)} implies \SGARefLink{sga3:xiii:theorem:5:1:condition:iv}{(iv)}.
Conversely, if \SGARefLink{sga3:xiii:theorem:5:1:condition:iv}{(iv)} holds, then \(\mathfrak h\) contains a regular
element of \(\mathfrak g\), and since it has the same rank as
\(\mathfrak g\), it contains a Cartan subalgebra by \SGARefLink{sga3:xiii:proposition:4:6}{4.6}.

We next prove the equivalence of conditions \SGARefLink{sga3:xiii:theorem:5:1:condition:viii}{(viii)} through \SGARefLink{sga3:xiii:theorem:5:1:condition:x}{(x)}. We
first note \SGARefLink{sga3:xiii:lemma:5:2}{the following facts}.

\begin{lemma}
\SGARefTarget{sga3:xiii:lemma:5:2}
\begin{enumerate}
  \item[a)]\SGARefTarget{sga3:xiii:lemma:5:2:part:a}
  If \(N\) is smooth, then
  \[
    \dim X\leq\dim G,
  \]
  with equality if and only if \(\mathfrak h=\mathfrak n\).

  \item[b)]\SGARefTarget{sga3:xiii:lemma:5:2:part:b}
  If \(\mathfrak h=\mathfrak n\), then
  \[
    \dim X\geq\dim G,
  \]
  with equality if and only if \(N\) is smooth.
\end{enumerate}
\end{lemma}

These assertions follow at once from \SGARefLink{sga3:xiii:lemma:5:2:display:dimension-formula}{the formula}
\SGARefTarget{sga3:xiii:lemma:5:2:display:dimension-formula}
\[
  \dim X=\dim G-\dim N+\dim_k\mathfrak h
\]
and from the fact that
\(\dim N=\dim_k\mathfrak n\) is equivalent to the smoothness of \(N\).

Now \SGARefLink{sga3:xiii:theorem:5:1:condition:viii}{(viii)} implies \SGARefLink{sga3:xiii:theorem:5:1:condition:x}{(x)}: by dominance, \SGARefLink{sga3:xiii:theorem:5:1:condition:viii}{(viii)} gives
\(\dim X\geq\dim G\), so \SGARefLink{sga3:xiii:lemma:5:2:part:a}{Lemma~5.2(a)} gives equality and
\(\mathfrak h=\mathfrak n\), hence \SGARefLink{sga3:xiii:theorem:5:1:condition:x}{(x)}. The same argument in reverse,
using \SGARefLink{sga3:xiii:lemma:5:2:part:b}{Lemma~5.2(b)}, shows that \SGARefLink{sga3:xiii:theorem:5:1:condition:x}{(x)} implies \SGARefLink{sga3:xiii:theorem:5:1:condition:viii}{(viii)}. Moreover, \SGARefLink{sga3:xiii:theorem:5:1:condition:viii}{(viii)}
implies \SGARefLink{sga3:xiii:theorem:5:1:condition:ix}{(ix)}, since it implies \(\mathfrak h=\mathfrak n\), so
\(\mathfrak h\) is the Lie algebra of the smooth algebraic subgroup
\(N\) of \(G\). Conversely, \SGARefLink{sga3:xiii:theorem:5:1:condition:ix}{(ix)} implies \SGARefLink{sga3:xiii:theorem:5:1:condition:viii}{(viii)}: since \(H\)
normalizes its Lie algebra \(\mathfrak h\), it is contained in \(N\);
and since \(H\) is smooth and has the same Lie algebra as \(N\), it
follows that \(N\) is smooth.

Finally, we prove the equivalence of \SGARefLink{sga3:xiii:theorem:5:1:condition:xi}{(xi)}, \SGARefLink{sga3:xiii:theorem:5:1:condition:xii}{(xii)}, and \SGARefLink{sga3:xiii:theorem:5:1:condition:xiii}{(xiii)}, and show
that they imply \SGARefLink{sga3:xiii:theorem:5:1:condition:iii}{(iii)}; this will complete \SGARefLink{sga3:xiii:theorem:5:1:diagram:implications}{the implication diagram}.
We have \SGARefLink{sga3:xiii:theorem:5:1:condition:xi}{(xi)}\(\Longleftrightarrow\)\SGARefLink{sga3:xiii:theorem:5:1:condition:xiii}{(xiii)}. Indeed, if
\(\mathfrak n=\mathfrak h\), then \SGARefLink{sga3:xiii:lemma:5:2:part:b}{Lemma~5.2(b)} gives
\(\dim X\geq\dim G\). Since \(W(\mathfrak g)\) is normal, condition
\SGARefLink{sga3:xiii:theorem:5:1:condition:xi}{(xi)} is therefore equivalent to \(\psi\) being generically
unramified, and this in turn is equivalent to \SGARefLink{sga3:xiii:theorem:5:1:condition:xiii}{(xiii)} by
\SGARefLink{sga3:xiii:lemma:1:1}{1.1}, \SGARefLink{sga3:xiii:theorem:5:1:condition:ii}{(ii)}\(\Longleftrightarrow\)\SGARefLink{sga3:xiii:theorem:5:1:condition:iii}{(iii)}, arguing as above in the proof
of \SGARefLink{sga3:xiii:theorem:5:1:condition:ii}{(ii)}\(\Longleftrightarrow\)\SGARefLink{sga3:xiii:theorem:5:1:condition:iii}{(iii)}.

Since \SGARefLink{sga3:xiii:theorem:5:1:condition:xi}{(xi)}\(\Rightarrow\)\SGARefLink{sga3:xiii:theorem:5:1:condition:x}{(x)}\(\Rightarrow\)\SGARefLink{sga3:xiii:theorem:5:1:condition:viii}{(viii)}, what we have
already proved shows that \SGARefLink{sga3:xiii:theorem:5:1:condition:xi}{(xi)} implies that \(N\) is smooth. Thus
\(q:G\times W(\mathfrak h)\to X\) is smooth, and the composite
\(\varphi=\psi\circ q\) is generically smooth; this is \SGARefLink{sga3:xiii:theorem:5:1:condition:iii}{(iii)}. Since
\SGARefLink{sga3:xiii:theorem:5:1:condition:iii}{(iii)} implies \SGARefLink{sga3:xiii:theorem:5:1:condition:i}{(i)}, it also follows that
\SGARefLink{sga3:xiii:theorem:5:1:condition:xi}{(xi)}\(\Rightarrow\)\SGARefLink{sga3:xiii:theorem:5:1:condition:xii}{(xii)}. Finally, \SGARefLink{sga3:xiii:theorem:5:1:condition:xii}{(xii)} implies \SGARefLink{sga3:xiii:theorem:5:1:condition:xi}{(xi)}. Indeed, \SGARefLink{sga3:xiii:theorem:5:1:condition:xii}{(xii)} plainly implies
\SGARefLink{sga3:xiii:theorem:5:1:condition:i}{(i)}, and we have seen that
\SGARefLink{sga3:xiii:theorem:5:1:condition:i}{(i)}\(\Rightarrow\)\SGARefLink{sga3:xiii:theorem:5:1:condition:iii}{(iii)}\(\Rightarrow\)\SGARefLink{sga3:xiii:theorem:5:1:condition:v}{(v)}; hence \SGARefLink{sga3:xiii:theorem:5:1:condition:xii}{(xii)} implies \SGARefLink{sga3:xiii:theorem:5:1:condition:v}{(v)}.
It also implies \SGARefLink{sga3:xiii:theorem:5:1:condition:ix}{(ix)}, and \SGARefLink{sga3:xiii:theorem:5:1:condition:ix}{(ix)} implies \SGARefLink{sga3:xiii:theorem:5:1:condition:x}{(x)}. Thus \SGARefLink{sga3:xiii:theorem:5:1:condition:xii}{(xii)} implies both
\SGARefLink{sga3:xiii:theorem:5:1:condition:v}{(v)} and \SGARefLink{sga3:xiii:theorem:5:1:condition:x}{(x)}, and therefore \SGARefLink{sga3:xiii:theorem:5:1:condition:xi}{(xi)}, since
\[
  \text{generically étale}
  =
  \text{generically smooth}
  +
  \text{generically quasi-finite}.
\]

If \(k\) has characteristic zero, then \SGARefLink{sga3:xiii:theorem:5:1:condition:vii}{(vii)} implies \SGARefLink{sga3:xiii:theorem:5:1:condition:viii}{(viii)}, because
Cartier's theorem makes \(N\) automatically smooth
(VI\textsubscript{B}~1.6.1). Moreover, [\SGARefLink{sga3:xiii:theorem:5:1:condition:viii}{(viii)} and \SGARefLink{sga3:xiii:theorem:5:1:condition:x}{(x)}] implies \SGARefLink{sga3:xiii:theorem:5:1:condition:xi}{(xi)},
because in characteristic zero, for a morphism of integral
preschemes,
\[
  \text{generically étale}
  =
  \text{dominant and generically quasi-finite}.
\]
Thus in this case all conditions \SGARefLink{sga3:xiii:theorem:5:1:condition:i}{(i)} through \SGARefLink{sga3:xiii:theorem:5:1:condition:xiii}{(xiii)} are equivalent.

\begin{corollary}
\SGARefTarget{sga3:xiii:corollary:5:3}
Under the equivalent conditions \SGARefLink{sga3:xiii:theorem:5:1:condition:viii}{(viii)} through \SGARefLink{sga3:xiii:theorem:5:1:condition:x}{(x)}, there exists a
unique smooth connected algebraic subgroup \(H\) of \(G\) whose Lie
algebra is \(\mathfrak h\), and
\[
  \operatorname{Norm}_G(H)
  =\operatorname{Norm}_G(\mathfrak h)
  =N,
  \qquad
  H=N^0.
\]
\end{corollary}

Indeed, \(H=N^0\) has the required properties. Conversely, if \(H\)
has these properties, then \(H\subset N\), since \(H\) normalizes its
Lie algebra \(\mathfrak h\). This is an inclusion of smooth groups
with the same Lie algebra, and \(H\) is connected, so \(H=N^0\). To
prove
\(\operatorname{Norm}_G(H)=\operatorname{Norm}_G(\mathfrak h)\), we
may assume that \(k\) is algebraically closed. What we have just proved
then shows immediately that the two groups have the same \(k\)-points.
On the other hand, the inclusions
\[
  H\subset\operatorname{Norm}_G(H)\subset N
\]
show that \(\operatorname{Norm}_G(H)\) and \(N\) have the same Lie
algebra; hence they are equal.

\begin{corollary}
\SGARefTarget{sga3:xiii:corollary:5:4}
Under the equivalent conditions \SGARefLink{sga3:xiii:theorem:5:1:condition:i}{(i)} through \SGARefLink{sga3:xiii:theorem:5:1:condition:iv}{(iv)}, let
\(a\in\mathfrak h\). Then the following conditions are equivalent, and
they hold whenever \(a\) is regular in \(\mathfrak g\):
\begin{enumerate}
  \item[(i)]\SGARefTarget{sga3:xiii:corollary:5:4:condition:i}
  \(\varphi\) is smooth at \((e,a)\).

  \item[(ii)]\SGARefTarget{sga3:xiii:corollary:5:4:condition:ii}
  \(M_a=\operatorname{Transp}_G(a,\mathfrak h)\) is smooth at \(e\),
  and
  \(\dim_e(M_a)=\operatorname{rank}_k\mathfrak h\).

  \item[(iii)]\SGARefTarget{sga3:xiii:corollary:5:4:condition:iii}
  \(\operatorname{ad}(a)_{\mathfrak g/\mathfrak h}\) is injective
  (equivalently, bijective).
\end{enumerate}
Under the equivalent conditions \SGARefLink{sga3:xiii:theorem:5:1:condition:xi}{(xi)} through \SGARefLink{sga3:xiii:theorem:5:1:condition:xiii}{(xiii)}, let \(H\) be the
algebraic subgroup of \(G\) considered in \SGARefLink{sga3:xiii:corollary:5:3}{Corollary~5.3}. The preceding
conditions are also equivalent to:
\begin{enumerate}
  \item[(iv)]\SGARefTarget{sga3:xiii:corollary:5:4:condition:iv}
  \(\psi\) is étale at \((\overline e,a)\).

  \item[(v)]\SGARefTarget{sga3:xiii:corollary:5:4:condition:v}
  If \(M_a^0\) denotes the connected component of \(e\) in the
  transporter \(M_a\) of \(a\) into \(\mathfrak h\), endowed with the
  structure induced by \(M_a\), then
  \[
    H=M_a^0.
  \]
\end{enumerate}
\end{corollary}

Plainly \SGARefLink{sga3:xiii:corollary:5:4:condition:i}{(i)} implies \SGARefLink{sga3:xiii:corollary:5:4:condition:ii}{(ii)}, since \(M_a\) is isomorphic to the fiber
\(\varphi^{-1}(a)\), with \(e\) corresponding to \((e,a)\).
Conversely, \SGARefLink{sga3:xiii:corollary:5:4:condition:ii}{(ii)} implies \SGARefLink{sga3:xiii:corollary:5:4:condition:i}{(i)}: condition \SGARefLink{sga3:xiii:corollary:5:4:condition:ii}{(ii)} says that \(\varphi\) is
``equidimensional'' at \((e,a)\), in the sense that the dimension of
the fiber through this point is the dimension of the generic fiber.
Since \(G\times W(\mathfrak h)\) and \(W(\mathfrak g)\) are regular,
this implies that \(\varphi\) is flat at \((e,a)\), hence smooth there
because its fiber is smooth at that point. The equivalence of \SGARefLink{sga3:xiii:corollary:5:4:condition:i}{(i)} and
\SGARefLink{sga3:xiii:corollary:5:4:condition:iii}{(iii)} was obtained in the proof of \SGARefLink{sga3:xiii:theorem:5:1}{Theorem~5.1} by the elementary
calculation of the tangent map. We also saw in \SGARefLink{sga3:xiii:proposition:4:8:item:b}{4.8(b)} that regularity
of \(a\) in \(\mathfrak g\) implies \SGARefLink{sga3:xiii:corollary:5:4:condition:iii}{(iii)}.

Under conditions \SGARefLink{sga3:xiii:theorem:5:1:condition:xi}{(xi)} through \SGARefLink{sga3:xiii:theorem:5:1:condition:xiii}{(xiii)}, the morphism
\(q:G\times W(\mathfrak h)\to X\) is smooth, since \(N\) is smooth.
Thus \SGARefLink{sga3:xiii:corollary:5:4:condition:i}{(i)} is equivalent to the smoothness of \(\psi\) at
\((\overline e,a)\); and since \(\psi\) is generically étale, this is
equivalent to \SGARefLink{sga3:xiii:corollary:5:4:condition:iv}{(iv)}. Finally, as was noted at the end of the proof of
\SGARefLink{sga3:xiii:lemma:1:1}{1.1}, condition \SGARefLink{sga3:xiii:corollary:5:4:condition:iv}{(iv)} implies that \(N\) is the prescheme induced on
\(M_a\) by an open and closed subset of \(M_a\), which gives \SGARefLink{sga3:xiii:corollary:5:4:condition:v}{(v)}.
Conversely, \SGARefLink{sga3:xiii:corollary:5:4:condition:v}{(v)} trivially implies \SGARefLink{sga3:xiii:corollary:5:4:condition:ii}{(ii)}. Alternatively, \SGARefLink{sga3:xiii:corollary:5:4:condition:v}{(v)} implies
\SGARefLink{sga3:xiii:corollary:5:4:condition:iv}{(iv)} by \SGARefLink{sga3:xiii:lemma:1:1}{1.1}, since \(\psi\) is dominant and \(W(\mathfrak g)\) is
normal, and ``unramified'' is equivalent here to ``étale''. This
completes the proof of \SGARefLink{sga3:xiii:corollary:5:4}{Corollary~5.4}.

\begin{corollary}
\SGARefTarget{sga3:xiii:corollary:5:5}
Let \(G\) be a smooth algebraic group over a field \(k\), and let
\(H\) be a smooth algebraic subgroup such that its Lie algebra
\(\mathfrak h\) contains, at least after a suitable extension of
\(k\), a Cartan subalgebra of the Lie algebra \(\mathfrak g\) of
\(G\). Let \(K\) be a connected algebraic subgroup of \(G\), not
necessarily smooth, with Lie algebra \(\mathfrak k\), and suppose that
\(\mathfrak k\) contains a regular element \(a\) of \(\mathfrak g\),
again at least after a suitable extension of \(k\). Then \(H\)
contains \(K\) if and only if \(\mathfrak h\) contains
\(\mathfrak k\).
\end{corollary}

Indeed, \SGARefLink{sga3:xiii:corollary:5:4}{Corollary~5.4}, \SGARefLink{sga3:xiii:corollary:5:4:condition:iii}{(iii)}\(\Rightarrow\)\SGARefLink{sga3:xiii:corollary:5:4:condition:v}{(v)}, applied to the
connected subgroup supplied by \SGARefLink{sga3:xiii:corollary:5:3}{Corollary~5.3}, gives
\(H^0=M_a^0\).\footnote{The French source prints \(H=M_a^0\) here. In the stated
generality the argument gives \(H^0=M_a^0\), which is all that is
needed.}
This identity is invariant under extension of the base field, so it
remains valid without assuming that the base field is infinite.
Moreover, \(\mathfrak k\subset\mathfrak h\) plainly implies
\(K\subset M_a\). Since \(K\) is connected,
\(K\subset M_a^0\subset H\). \hfill\(\square\)

\begin{corollary}
\SGARefTarget{sga3:xiii:corollary:5:6}
Let \(G\) and \(H\) be as in \SGARefLink{sga3:xiii:corollary:5:5}{Corollary~5.5}, and let \(K\) be an
algebraic subgroup of \(G\). Suppose that \(H\) is connected and \(K\)
is smooth. Then \(K\) contains \(H\) if and only if
\(\mathfrak k\) contains \(\mathfrak h\).
\end{corollary}

Indeed, if \(\mathfrak k\supset\mathfrak h\), then \(K\) satisfies,
with respect to \(H\), the hypothesis imposed on \(H\) in
\SGARefLink{sga3:xiii:corollary:5:5}{Corollary~5.5}. Meanwhile \(H\) plainly satisfies, with respect to
\(K\), the condition imposed on \(K\) there. The conclusion follows
from \SGARefLink{sga3:xiii:corollary:5:5}{Corollary~5.5}.

\begin{corollary}
\SGARefTarget{sga3:xiii:corollary:5:7}
Let \(\mathfrak g\) be the Lie algebra of a smooth algebraic group
\(G\) over a field \(k\). Then:
\begin{enumerate}
  \item[a)]\SGARefTarget{sga3:xiii:corollary:5:7:part:a}
  A Lie subalgebra \(\mathfrak d\) of \(\mathfrak g\) is a Cartan
  subalgebra if and only if it is nilpotent and equal to its own
  normalizer.

  \item[b)]\SGARefTarget{sga3:xiii:corollary:5:7:part:b}
  An element \(a\in\mathfrak g\) is regular if and only if
  \(\operatorname{Nil}(a,\mathfrak g)\) is nilpotent.
\end{enumerate}
\end{corollary}

After extending the base field, we may assume that \(k\) is infinite.
In view of \SGARefLink{sga3:xiii:proposition:4:4}{4.4}, to prove \SGARefLink{sga3:xiii:corollary:5:7:part:a}{(a)} it remains only to show that if
\(\mathfrak d\) is nilpotent and contains an element \(a\) for which
\(\operatorname{ad}(a)_{\mathfrak g/\mathfrak d}\) is injective, then
\(\mathfrak d\) is a Cartan subalgebra. By \SGARefLink{sga3:xiii:theorem:5:1}{Theorem~5.1},
\SGARefLink{sga3:xiii:theorem:5:1:condition:ii}{(ii)}\(\Rightarrow\)\SGARefLink{sga3:xiii:theorem:5:1:condition:i}{(i)}, \(\mathfrak d\) contains a Cartan subalgebra
\[
  \mathfrak d'=\operatorname{Nil}(a,\mathfrak g).
\]
Since \(\mathfrak d\) is nilpotent, \SGARefLink{sga3:xiii:corollary:4:3}{4.3} now gives
\(\mathfrak d=\mathfrak d'\). For \SGARefLink{sga3:xiii:corollary:5:7:part:b}{(b)}, one notes that
\(\operatorname{Nil}(a,\mathfrak g)\) is a Cartan subalgebra of
\(\mathfrak g\) by \SGARefLink{sga3:xiii:corollary:5:7:part:a}{(a)}, and therefore \(a\) is regular.

\begin{corollary}
\SGARefTarget{sga3:xiii:corollary:5:8}
Let \(G\) be a smooth algebraic group over a field \(k\), and let
\(H\) be a smooth algebraic subgroup, with Lie algebras
\(\mathfrak g\) and \(\mathfrak h\), respectively. Suppose that after
a suitable extension of the base field, \(\mathfrak h\) contains a
Cartan subalgebra of \(\mathfrak g\). Let \(\mathfrak d\) be a Lie
subalgebra of \(\mathfrak h\). Then \(\mathfrak d\) is a Cartan
subalgebra of \(\mathfrak h\) if and only if it is a Cartan subalgebra
of \(\mathfrak g\).
\end{corollary}

In view of \SGARefLink{sga3:xiii:proposition:4:8:item:c}{4.8(c)}, it remains to show that if \(\mathfrak d\) is a
Cartan subalgebra of \(\mathfrak h\), then it is one of
\(\mathfrak g\). It is enough to show that \(\mathfrak d\) contains
an element \(a\) regular in \(\mathfrak g\), and we may assume that
\(k\) is algebraically closed. But there is a dense open subset of
\(\mathfrak h\) consisting of points regular in \(\mathfrak g\), so
the assertion follows from \SGARefLink{sga3:xiii:theorem:5:1}{Theorem~5.1},
\SGARefLink{sga3:xiii:theorem:5:1:condition:i}{(i)}\(\Rightarrow\)\SGARefLink{sga3:xiii:theorem:5:1:condition:vii}{(vii)}, applied to
\((\mathfrak h,\mathfrak d)\).

\section{Cartan subalgebras and subgroups of type \texorpdfstring{\((C)\)}{(C)}
associated with a smooth algebraic group}
\SGARefTarget{sga3:xiii:section:6}

For simplicity, in \SGARefLink{sga3:xiii:theorem:6:1}{the following theorem} we restrict ourselves to an
algebraically closed base field; the case of an arbitrary base
prescheme will be treated in the next Exposé.

\begin{theorem}
\SGARefTarget{sga3:xiii:theorem:6:1}
Let \(G\) be a smooth algebraic group over an algebraically closed
field \(k\), and let \(\mathfrak g\) be its Lie algebra. Then:
\begin{enumerate}
  \item[a)]\SGARefTarget{sga3:xiii:theorem:6:1:part:a}
  The Cartan subalgebras of \(\mathfrak g\) are conjugate.

  \item[b)]\SGARefTarget{sga3:xiii:theorem:6:1:part:b}
  Let \(\mathfrak d\) be a Cartan subalgebra of \(\mathfrak g\).
  Its normalizer \(N\) in \(G\) is smooth, and \(D=N^0\) is the
  unique smooth connected subgroup of \(G\) whose Lie algebra is
  \(\mathfrak d\). One has
  \[
    \operatorname{Norm}_G(\mathfrak d)
    =\operatorname{Norm}_G(D)
    =N,
    \qquad\text{and hence}\qquad
    D=\operatorname{Norm}_G(D)^0.
  \]

  \item[c)]\SGARefTarget{sga3:xiii:theorem:6:1:part:c}
  With \(\mathfrak d\) as in \SGARefLink{sga3:xiii:theorem:6:1:part:b}{(b)}, set, as in \SGARefLink{sga3:xiii:section:5}{Section~5},
  \[
    X=G\times^N W(\mathfrak d),
  \]
  and consider the canonical morphism
  \[
    \psi:X\longrightarrow W(\mathfrak g).
  \]
  Its fiber over \(a\in\mathfrak g\) has as its \(k\)-valued points
  the set of Cartan subalgebras of \(\mathfrak g\) containing \(a\).
  Then \(\psi\) is birational.

  \item[d)]\SGARefTarget{sga3:xiii:theorem:6:1:part:d}
  With the notation of \SGARefLink{sga3:xiii:theorem:6:1:part:c}{(c)}, let \(U\) be the largest open subset of
  \(W(\mathfrak g)\) for which \(\psi\) induces an isomorphism
  \[
    \psi^{-1}(U)\xrightarrow{\sim}U.
  \]
  Then, for \(a\in\mathfrak g\), the following conditions are
  equivalent:
  \begin{enumerate}
    \item[(i)]\SGARefTarget{sga3:xiii:theorem:6:1:part:d:condition:i}
    \(a\in U(k)\).

    \item[(ii)]\SGARefTarget{sga3:xiii:theorem:6:1:part:d:condition:ii}
    The element \(a\) belongs to one and only one Cartan subalgebra of
    \(\mathfrak g\).

    \item[(iii)]\SGARefTarget{sga3:xiii:theorem:6:1:part:d:condition:iii}
    The set of Cartan subalgebras of \(\mathfrak g\) containing \(a\)
    is finite and nonempty.

    \item[(iv)]\SGARefTarget{sga3:xiii:theorem:6:1:part:d:condition:iv}
    The fiber \(\psi^{-1}(a)\) has an isolated point.

    \item[(v)]\SGARefTarget{sga3:xiii:theorem:6:1:part:d:condition:v}
    If \(a\in\mathfrak d\), the morphism \(\psi\) is étale---or merely
    quasi-finite---at the point \((\overline e,a)\).

    \item[(vi)]\SGARefTarget{sga3:xiii:theorem:6:1:part:d:condition:vi}
    The element \(a\) is regular in \(\mathfrak g\).

    \item[(vii)]\SGARefTarget{sga3:xiii:theorem:6:1:part:d:condition:vii}
    With the notation of \SGARefLink{sga3:xiii:corollary:5:4:condition:iii}{Corollary~5.4(iii)}, one has \(N=M_a\).
  \end{enumerate}
\end{enumerate}
\end{theorem}

\begin{proof}
Apply \SGARefLink{sga3:xiii:theorem:5:1}{Theorem~5.1} with the Cartan subalgebra \(\mathfrak d\) of
\(\mathfrak g\) in place of \(\mathfrak h\). We show that the strongest
conditions, \SGARefLink{sga3:xiii:theorem:5:1:condition:xi}{(xi)} through \SGARefLink{sga3:xiii:theorem:5:1:condition:xiii}{(xiii)}, hold. This is clear either in the
form \SGARefLink{sga3:xiii:theorem:5:1:condition:xiii}{(xiii)}, by \SGARefLink{sga3:xiii:corollary:4:7}{4.7}---which says that if a regular element
\(a\in\mathfrak g\) lies in \(\mathfrak d\), then \(N=M_a\)---or in
the form
\[
  \SGARefLink{sga3:xiii:theorem:5:1:condition:xi}{(\mathrm{xi})}=\SGARefLink{sga3:xiii:theorem:5:1:condition:x}{(\mathrm{x})}+\SGARefLink{sga3:xiii:theorem:5:1:condition:i}{(\mathrm{i})}.
\]
Indeed, \SGARefLink{sga3:xiii:theorem:5:1:condition:i}{(i)} is automatic, while \SGARefLink{sga3:xiii:theorem:5:1:condition:x}{(x)} follows from the fact that a
regular point of \(\mathfrak g\) belongs to a unique Cartan subalgebra
of \(\mathfrak g\) (\SGARefLink{sga3:xiii:proposition:4:6:item:a}{4.6(a)}), and therefore, a fortiori, to a unique
conjugate of \(\mathfrak d\).

Part \SGARefLink{sga3:xiii:theorem:6:1:part:b}{(b)} now follows from \SGARefLink{sga3:xiii:corollary:5:3}{Corollary~5.3}. Part \SGARefLink{sga3:xiii:theorem:6:1:part:c}{(c)} follows because
\(\psi\) is generically étale and because a sufficiently general
point of \(\mathfrak g\)---more precisely, a regular point---belongs
to a unique Cartan subalgebra of \(\mathfrak g\). Under these
conditions, the equivalence of conditions \SGARefLink{sga3:xiii:theorem:6:1:part:d:condition:i}{(i)} through \SGARefLink{sga3:xiii:theorem:6:1:part:d:condition:v}{(v)} on \(a\)
follows immediately from Zariski's Main Theorem, applied to the
separated birational morphism \(\psi\), since \(W(\mathfrak g)\) is
normal and \(X\) is integral. The equivalence of \SGARefLink{sga3:xiii:theorem:6:1:part:d:condition:v}{(v)} and \SGARefLink{sga3:xiii:theorem:6:1:part:d:condition:vi}{(vi)} is the
special case of \SGARefLink{sga3:xiii:corollary:5:4}{Corollary~5.4},
\SGARefLink{sga3:xiii:theorem:6:1:part:d:condition:iii}{(iii)}\(\Longleftrightarrow\)\SGARefLink{sga3:xiii:theorem:6:1:part:d:condition:iv}{(iv)}, obtained after carrying \(a\) into
\(\mathfrak d\) by a suitable element \(g\in G(k)\), in view of
\SGARefLink{sga3:xiii:proposition:4:6:item:b}{4.6(b)}. Moreover, \SGARefLink{sga3:xiii:corollary:5:4}{Corollary~5.4} shows that \SGARefLink{sga3:xiii:theorem:6:1:part:d:condition:v}{(v)} and \SGARefLink{sga3:xiii:theorem:6:1:part:d:condition:vi}{(vi)} are also
equivalent to
\[
  N\supset M_a^0.
\]
By \SGARefLink{sga3:xiii:corollary:4:7}{4.7} this even implies
\(N=M_a\).\footnote{The French source prints \(N=M\) at this point; condition
\SGARefLink{sga3:xiii:theorem:6:1:part:d:condition:vii}{(vii)} and the surrounding notation require \(N=M_a\).}
This proves \SGARefLink{sga3:xiii:theorem:6:1:part:d}{(d)}.

Of course, \SGARefLink{sga3:xiii:theorem:6:1:part:b}{(b)}, \SGARefLink{sga3:xiii:theorem:6:1:part:c}{(c)}, and the equivalence of \SGARefLink{sga3:xiii:theorem:6:1:part:d:condition:i}{(i)}, \SGARefLink{sga3:xiii:theorem:6:1:part:d:condition:iv}{(iv)}, \SGARefLink{sga3:xiii:theorem:6:1:part:d:condition:v}{(v)}, \SGARefLink{sga3:xiii:theorem:6:1:part:d:condition:vi}{(vi)}, and
\SGARefLink{sga3:xiii:theorem:6:1:part:d:condition:vii}{(vii)} remain valid over an arbitrary base field. We now prove \SGARefLink{sga3:xiii:theorem:6:1:part:a}{(a)},
using the assumption that \(k\) is algebraically closed. By
\SGARefLink{sga3:xiii:theorem:5:1}{Theorem~5.1}, \SGARefLink{sga3:xiii:theorem:5:1:condition:i}{(i)}\(\Rightarrow\)\SGARefLink{sga3:xiii:theorem:5:1:condition:vii}{(vii)}, the morphism
\(\psi:X\to W(\mathfrak g)\) is dominant. Hence there is a dense open
subset \(V\) of \(W(\mathfrak g)\) such that every \(a\in V(k)\) is
the image of a point of \(X(k)\), or equivalently, belongs to a
conjugate of \(\mathfrak d\). Apply \SGARefLink{sga3:xiii:theorem:5:1:diagram:implications}{the same result} to a second Cartan
subalgebra \(\mathfrak d'\) of \(\mathfrak g\). We may choose \(V\)
so that every element of \(V(k)\) is conjugate both to an element of
\(\mathfrak d\) and to an element of \(\mathfrak d'\). Choose a
regular element of \(V(k)\). Some conjugate \(\mathfrak d''\) of
\(\mathfrak d'\) then contains an element of \(\mathfrak d\) regular
in \(\mathfrak g\), and \SGARefLink{sga3:xiii:proposition:4:6:item:a}{4.6(a)} gives
\(\mathfrak d''=\mathfrak d\). This proves \SGARefLink{sga3:xiii:theorem:6:1:part:a}{(a)} and completes the proof.
\end{proof}

\begin{definition}
\SGARefTarget{sga3:xiii:definition:6:2}
Let \(G\) be a smooth algebraic group over a field \(k\). A
\emph{subgroup of type \((C)\)} of \(G\) is a smooth connected
algebraic subgroup whose Lie algebra is a Cartan subalgebra of
\(\mathfrak g\). The \emph{infinitesimal rank} of \(G\) is the
nilpotent rank of its Lie algebra \(\mathfrak g\); it is equal to the
dimension of every subgroup of type \((C)\) of \(G\).
\end{definition}

\SGARefLink{sga3:xiii:theorem:6:1:part:b}{Theorem~6.1(b)} immediately gives:

\begin{corollary}
\SGARefTarget{sga3:xiii:corollary:6:3}
Under the hypotheses of \SGARefLink{sga3:xiii:definition:6:2}{Definition~6.2}, the map
\[
  D\longmapsto\mathfrak d=\operatorname{Lie}(D)
\]
establishes a bijection between the subgroups of type \((C)\) of
\(G\) and the Cartan subalgebras \(\mathfrak d\) of \(\mathfrak g\).
If \(D\) is a subgroup of type \((C)\) of \(G\), then \(D\) is its
own connected normalizer:
\[
  D=\operatorname{Norm}_G(D)^0.
\]
\end{corollary}

Combining \SGARefLink{sga3:xiii:theorem:6:1:part:a}{Theorem~6.1(a)} and \SGARefLink{sga3:xiii:corollary:6:3}{Corollary~6.3} gives:

\begin{corollary}
\SGARefTarget{sga3:xiii:corollary:6:4}
If \(k\) is algebraically closed, the subgroups of type \((C)\) of
\(G\) are conjugate to one another.
\end{corollary}

\begin{corollary}
\SGARefTarget{sga3:xiii:corollary:6:5}
Let \(G\) be a smooth algebraic group over an algebraically closed
field \(k\), and let \(H\) be a smooth algebraic subgroup of \(G\),
with Lie algebras \(\mathfrak g\) and \(\mathfrak h\), respectively.
Then \(\mathfrak h\) contains a Cartan subalgebra of \(\mathfrak g\)
if and only if \(H\) contains a subgroup \(D\) of type \((C)\) of
\(G\).
\end{corollary}

The condition is plainly sufficient. It is also necessary, since by
\SGARefLink{sga3:xiii:corollary:5:5}{Corollary~5.5},
\[
  H\supset D
  \quad\Longleftrightarrow\quad
  \mathfrak h\supset\mathfrak d.
\]

\refstepcounter{theorem}
\par\medskip
\noindent\SGARefTarget{sga3:xiii:remarks:6:6}%
\textbf{Remarks \thetheorem.}
\begin{enumerate}
  \item[a)]\SGARefTarget{sga3:xiii:remarks:6:6:part:a}
  The connected subgroups \(H\) of \(G\) described in
  \SGARefLink{sga3:xiii:corollary:6:5}{Corollary~6.5} correspond bijectively to the Lie subalgebras
  \(\mathfrak h\) of \(\mathfrak g\) satisfying the strongest
  condition, \SGARefLink{sga3:xiii:theorem:5:1:condition:xii}{(xii)}, of \SGARefLink{sga3:xiii:theorem:5:1}{Theorem~5.1}. By \SGARefLink{sga3:xiii:corollary:5:5}{Corollary~5.5}, inclusion
  relations among such subgroups can already be read from their Lie
  algebras.

  \item[b)]\SGARefTarget{sga3:xiii:remarks:6:6:part:b}
  Continue to suppose that \(k\) is algebraically closed, and let
  \(D\) be a subgroup of type \((C)\) of \(G\). It is then easy to
  show that \(D\) contains a Cartan subgroup \(C\) of \(G\). Indeed,
  put \(V=\mathfrak g/\mathfrak d\). For a general element
  \(a\in\mathfrak d\), the endomorphism
  \(\operatorname{ad}(a)\) of \(V\) is injective; it suffices to take
  \(a\) regular in \(\mathfrak g\). From this one readily deduces that,
  for a general \(g\in D(k)\), the action of
  \(\operatorname{Ad}(g)\) on \(V\) has no nonzero invariant vector.
  Thus \SGARefLink{sga3:xiii:theorem:2:1}{2.1}, \SGARefLink{sga3:xiii:theorem:2:1:item:vii}{(vii)}\(\Rightarrow\)\SGARefLink{sga3:xiii:theorem:2:1:item:i}{(i)}, applies. In fact, in the next
  Exposé we shall prove the more precise result that every Cartan
  subgroup \(C\) of \(G\) is contained in one and only one subgroup
  of type \((C)\) of \(G\).

  \item[c)]\SGARefTarget{sga3:xiii:remarks:6:6:part:c}
  In general one must not confuse Cartan subgroups with subgroups of
  type \((C)\). The subgroups of type \((C)\) of \(G\) are Cartan
  subgroups if and only if they are nilpotent, since Cartan subgroups
  are maximal nilpotent subgroups. It may happen that \(\mathfrak g\)
  is nilpotent although \(G\) is not; an example is
  \(\operatorname{SL}(2)_k\) when \(k\) has characteristic \(2\).
  Then the Cartan subalgebras of \(\mathfrak g\) are all equal to
  \(\mathfrak g\). Thus, if \(G\) is connected, \(G\) is a subgroup
  of type \((C)\) of itself, but it is not a Cartan subgroup of
  itself.

  On the other hand, we shall see in Exposé~XIV that if \(G\) is an
  adjoint semisimple group, its subgroups of type \((C)\) are precisely
  its Cartan subgroups. The same conclusion holds in characteristic
  zero for every smooth algebraic group \(G\) over \(k\). This follows
  at once from the fact that, in characteristic zero, a connected
  algebraic group is nilpotent if and only if its Lie algebra is
  nilpotent. In turn, this follows from the fact that in
  characteristic zero the center \(Z\) of a connected algebraic group
  \(G\) has Lie algebra equal to the center of \(\mathfrak g\).
  Indeed, one obtains a priori the space of invariants in
  \(\mathfrak g\) under the adjoint representation of \(G\). Since
  \(G\) is connected and \(k\) has characteristic zero, these are also
  the infinitesimal invariants of the adjoint representation of
  \(\mathfrak g\), namely the center of \(\mathfrak g\). Finally,
  Cartier's theorem (cf.\ VI\textsubscript{B}~1.6.1) shows that \(Z\)
  is smooth.
\end{enumerate}
